\documentclass[12pt,a4paper]{article}

\usepackage[utf8]{inputenc}
\usepackage[T5]{fontenc}
\usepackage[vietnamese]{babel}
\usepackage{amsmath,amssymb}
\usepackage{array}
\usepackage{booktabs}
\usepackage{geometry}
\usepackage{enumitem}
\usepackage{fancyhdr}
\usepackage{hyperref}

\geometry{
    left=2.8cm,
    right=2.2cm,
    top=2.5cm,
    bottom=2.5cm
}

\hypersetup{
    unicode=true,
    colorlinks=true,
    linkcolor=black,
    citecolor=black,
    urlcolor=black
}

\pagestyle{fancy}
\fancyhf{}
\lhead{Cơ sở trí tuệ nhân tạo}
\rhead{Minimax cho Tic-Tac-Toe}
\cfoot{\thepage}

\setlength{\parindent}{0pt}
\setlength{\parskip}{6pt}

\newcommand{\blank}{\_}
\newcommand{\board}[9]{
\[
\begin{array}{c|c|c}
#1 & #2 & #3\\ \hline
#4 & #5 & #6\\ \hline
#7 & #8 & #9
\end{array}
\]
}

\begin{document}

\begin{titlepage}
    \centering
    {\Large \textbf{Bài tập nhóm giữa kỳ môn Cơ sở trí tuệ nhân tạo}\par}
    \vspace{1.2cm}

    {\large Giảng viên: TS. Nguyễn Bích Vân\par}
    \vspace{0.4cm}
    {\large Hạn nộp bài: 04/05/2026\par}

    \vspace{1.8cm}
    {\large \textbf{Đề tài: Giải bài toán Tic-Tac-Toe bằng thuật toán Minimax}\par}

    \vspace{1.5cm}
    \begin{tabular}{ll}
        \textbf{Nhóm 9} & \\
        Nguyễn Đức Trung & 24022473\\
        Phạm Sỹ Toàn & 24022467\\
        Lê Minh Đức & 24022291\\
        Phạm Công Thành & 24022455\\
        Trần Công Tuấn & 24022485
    \end{tabular}

    \vfill
    {\large 2026\par}
\end{titlepage}

\tableofcontents
\newpage

\section{Tóm tắt}

Báo cáo trình bày cách giải trò chơi Tic-Tac-Toe bằng thuật toán Minimax.
Trong trò chơi này, người chơi MAX là $X$, người chơi MIN là $O$.
Hàm đánh giá nhận giá trị $+1$ nếu $X$ thắng, $-1$ nếu $O$ thắng và $0$ nếu ván cờ hòa.
Mục tiêu của thuật toán là xác định nước đi tốt nhất tại mỗi trạng thái bàn cờ và tính giá trị Minimax tương ứng.

\section{Phát biểu bài toán}

Xét trò chơi Tic-Tac-Toe trên bàn cờ $3 \times 3$.
Hai người chơi lần lượt đánh dấu vào các ô trống.
Người chơi $X$ đi theo chiến lược cực đại hóa điểm số, còn người chơi $O$ đi theo chiến lược cực tiểu hóa điểm số.

Bài toán yêu cầu:
\begin{enumerate}[label=\arabic*.]
    \item Mô tả thuật toán Minimax cho phép người chơi đấu với máy.
    \item In ra nước đi tốt nhất và giá trị Minimax.
    \item Nhận xét về kết quả và đặc điểm của thuật toán.
\end{enumerate}

\section{Mô hình hóa trò chơi}

\subsection{Trạng thái bàn cờ}

Một trạng thái bàn cờ được ký hiệu là $s$.
Mỗi ô trên bàn cờ có thể nhận một trong ba giá trị:
\[
\{X, O, \blank\},
\]
trong đó $\blank$ biểu diễn ô trống.

Ví dụ, một trạng thái bàn cờ có thể được biểu diễn như sau:
\board{X}{O}{X}{\blank}{O}{\blank}{\blank}{X}{\blank}

\subsection{Tập nước đi hợp lệ}

Với một trạng thái $s$, ký hiệu $A(s)$ là tập tất cả các nước đi hợp lệ.
Một nước đi hợp lệ là một ô còn trống trên bàn cờ.
Nếu người chơi chọn nước đi $a \in A(s)$, trạng thái mới thu được được ký hiệu là:
\[
Result(s,a).
\]

\subsection{Trạng thái kết thúc}

Một trạng thái được gọi là trạng thái kết thúc nếu xảy ra một trong các trường hợp:
\begin{enumerate}[label=\arabic*.]
    \item Người chơi $X$ có ba quân liên tiếp trên một hàng, một cột hoặc một đường chéo.
    \item Người chơi $O$ có ba quân liên tiếp trên một hàng, một cột hoặc một đường chéo.
    \item Không còn ô trống nào trên bàn cờ và không người chơi nào thắng.
\end{enumerate}

Các đường thắng cần kiểm tra gồm:
\[
\begin{aligned}
&(1,1),(1,2),(1,3),\\
&(2,1),(2,2),(2,3),\\
&(3,1),(3,2),(3,3),\\
&(1,1),(2,1),(3,1),\\
&(1,2),(2,2),(3,2),\\
&(1,3),(2,3),(3,3),\\
&(1,1),(2,2),(3,3),\\
&(1,3),(2,2),(3,1).
\end{aligned}
\]

\section{Hàm đánh giá}

Hàm đánh giá trạng thái kết thúc được định nghĩa như sau:
\[
Eval(s)=
\begin{cases}
+1, & \text{nếu } X \text{ thắng},\\
-1, & \text{nếu } O \text{ thắng},\\
0, & \text{nếu ván cờ hòa}.
\end{cases}
\]

Vì $X$ là người chơi MAX nên $X$ luôn cố gắng chọn nước đi làm giá trị đánh giá lớn nhất.
Ngược lại, $O$ là người chơi MIN nên $O$ luôn cố gắng chọn nước đi làm giá trị đánh giá nhỏ nhất.

\section{Thuật toán Minimax}

Thuật toán Minimax xây dựng cây trò chơi từ trạng thái hiện tại.
Mỗi đỉnh của cây là một trạng thái bàn cờ.
Mỗi cạnh ứng với một nước đi hợp lệ.
Các nút lá là những trạng thái kết thúc và được gán giá trị bởi hàm $Eval(s)$.

Giá trị Minimax của một trạng thái $s$, ký hiệu $V(s)$, được định nghĩa đệ quy như sau:
\[
V(s)=
\begin{cases}
Eval(s), & \text{nếu } s \text{ là trạng thái kết thúc},\\
\max\limits_{a \in A(s)} V(Result(s,a)), & \text{nếu đến lượt } X,\\
\min\limits_{a \in A(s)} V(Result(s,a)), & \text{nếu đến lượt } O.
\end{cases}
\]

Ý nghĩa của công thức trên là:
\begin{itemize}
    \item Nếu trò chơi đã kết thúc, giá trị của trạng thái chính là giá trị đánh giá.
    \item Nếu đến lượt $X$, thuật toán chọn giá trị lớn nhất trong các trạng thái con.
    \item Nếu đến lượt $O$, thuật toán chọn giá trị nhỏ nhất trong các trạng thái con.
\end{itemize}

\subsection{Quy tắc chọn nước đi tốt nhất}

Nếu máy đóng vai $X$, nước đi tốt nhất tại trạng thái $s$ là:
\[
a^*=\arg\max_{a \in A(s)} V(Result(s,a)).
\]

Nếu máy đóng vai $O$, nước đi tốt nhất tại trạng thái $s$ là:
\[
a^*=\arg\min_{a \in A(s)} V(Result(s,a)).
\]

Trong đó $a^*$ là nước đi tối ưu mà máy cần chọn.

\subsection{Các bước thực hiện}

Quá trình tìm nước đi tốt nhất bằng Minimax có thể mô tả như sau:
\begin{enumerate}[label=\textbf{Bước \arabic*:}]
    \item Kiểm tra trạng thái hiện tại có phải là trạng thái kết thúc hay không.
    \item Nếu là trạng thái kết thúc, trả về giá trị $Eval(s)$.
    \item Nếu đến lượt $X$, xét toàn bộ nước đi hợp lệ và chọn giá trị lớn nhất.
    \item Nếu đến lượt $O$, xét toàn bộ nước đi hợp lệ và chọn giá trị nhỏ nhất.
    \item Trả về nước đi có giá trị Minimax tốt nhất đối với người chơi hiện tại.
\end{enumerate}

\section{Kết quả thực nghiệm}

\subsection{Trạng thái ban đầu}

Xét bàn cờ ban đầu rỗng:
\board{\blank}{\blank}{\blank}{\blank}{\blank}{\blank}{\blank}{\blank}{\blank}

Với trạng thái ban đầu $s_0$, nếu cả hai người chơi đều đi tối ưu thì trò chơi Tic-Tac-Toe sẽ kết thúc với kết quả hòa.
Do đó:
\[
V(s_0)=0.
\]

Một nước đi tối ưu cho $X$ là ô trung tâm:
\[
a^*=(2,2).
\]

Sau khi $X$ đi vào ô trung tâm, bàn cờ trở thành:
\board{\blank}{\blank}{\blank}{\blank}{X}{\blank}{\blank}{\blank}{\blank}

Giá trị Minimax của trạng thái sau nước đi này là:
\[
V(Result(s_0,a^*))=0.
\]

Điều này cho thấy $X$ không thể đảm bảo chiến thắng nếu $O$ cũng đi tối ưu, nhưng $X$ có thể đảm bảo không thua.

\subsection{Ví dụ trạng thái có nước thắng ngay cho MAX}

Xét trạng thái:
\board{X}{X}{\blank}{O}{O}{\blank}{\blank}{\blank}{\blank}

Đến lượt $X$, nước đi tốt nhất là đặt $X$ vào ô $(1,3)$:
\[
a^*=(1,3).
\]

Khi đó bàn cờ trở thành:
\board{X}{X}{X}{O}{O}{\blank}{\blank}{\blank}{\blank}

Vì $X$ thắng trên hàng thứ nhất nên:
\[
V=+1.
\]

\subsection{Ví dụ trạng thái cần chặn MIN}

Xét trạng thái:
\board{O}{O}{\blank}{X}{\blank}{\blank}{X}{\blank}{\blank}

Đến lượt $X$, nếu không chặn ô $(1,3)$ thì $O$ sẽ thắng ở lượt tiếp theo.
Vì vậy nước đi tốt nhất của $X$ là:
\[
a^*=(1,3).
\]

Sau nước đi này:
\board{O}{O}{X}{X}{\blank}{\blank}{X}{\blank}{\blank}

Nước đi này không tạo ra chiến thắng ngay cho $X$, nhưng ngăn $O$ đạt giá trị $-1$ ở lượt kế tiếp.
Theo quan điểm Minimax, đây là nước đi tối ưu vì nó tránh trạng thái thua chắc chắn.

\section{Đánh giá độ phức tạp}

Trong Tic-Tac-Toe, số ô tối đa là $9$.
Ở lượt đầu tiên có tối đa $9$ nước đi, lượt tiếp theo có tối đa $8$ nước đi, sau đó là $7$ nước đi, và cứ tiếp tục như vậy.
Do đó số trạng thái tối đa cần xét bị chặn bởi:
\[
9! = 362880.
\]

Về tổng quát, nếu $b$ là hệ số phân nhánh trung bình và $d$ là độ sâu tối đa của cây trò chơi, độ phức tạp thời gian của Minimax là:
\[
O(b^d).
\]

Độ phức tạp bộ nhớ phụ thuộc vào cách cài đặt.
Nếu dùng đệ quy theo chiều sâu, bộ nhớ cần dùng xấp xỉ:
\[
O(d).
\]

\section{Nhận xét}

Thuật toán Minimax phù hợp để giải Tic-Tac-Toe vì không gian trạng thái của trò chơi nhỏ.
Thuật toán có thể duyệt toàn bộ cây trò chơi để tìm ra nước đi tối ưu.
Nếu cả hai người chơi đều áp dụng chiến lược tối ưu, kết quả của trò chơi từ trạng thái ban đầu là hòa.

Ưu điểm chính của Minimax là cho lời giải chính xác và đảm bảo chọn được nước đi tốt nhất theo mô hình trò chơi đối kháng.
Tuy nhiên, nhược điểm của thuật toán là số trạng thái tăng rất nhanh khi kích thước trò chơi lớn hơn.
Với các trò chơi phức tạp như cờ vua hoặc cờ vây, việc duyệt toàn bộ cây trò chơi là không khả thi.
Khi đó cần dùng thêm các kỹ thuật như cắt tỉa Alpha--Beta, giới hạn độ sâu tìm kiếm hoặc hàm đánh giá xấp xỉ.

\section{Kết luận}

Bài toán Tic-Tac-Toe có thể được giải hiệu quả bằng thuật toán Minimax.
Bằng cách mô hình hóa $X$ là người chơi MAX và $O$ là người chơi MIN, thuật toán xét tất cả các nước đi hợp lệ, đánh giá các trạng thái kết thúc và lan truyền giá trị ngược lên cây trò chơi.
Kết quả cho thấy nước đi tối ưu từ bàn cờ ban đầu có giá trị Minimax bằng $0$, tức là hai bên đi tối ưu sẽ hòa.

\end{document}
